% Discussion, limitations and defenses (owner: C). Numbers from harness/, model/ and linking/
% hypotheses.md; the Track 3 parts are to be updated once the campaign results are final.

\section{Discussion}
\label{sec:discussion}

\paragraph{Agents look like their environment.} The central observation across our tracks is that a
coding agent's network identity is mostly not its own. When it shells out, the request is the system's
curl, or whatever Python the agent finds first on the path, byte for byte; the TLS and header layers
then describe the user's machine, not the agent. Only built-in tools such as Claude Code's WebFetch have
a fingerprint of their own, and that fingerprint is the same on every machine. This cuts both ways for
tracking. An agent that shells out is hard to tell from a script (H1) but carries the user's
environment, which links its sessions to each other and to the user's own command-line use; an agent
with its own client is easy to recognise but carries nothing about the user beyond the IP address.

\paragraph{What a passive site learns without any text.} Prompt injection is the obvious way for a site
to make an agent reveal more, but it is visible, detectable and increasingly resisted. Our results show
that a site learns a good deal from traffic it receives anyway. The pause before the submit separates an
LLM in the loop from a script; the sequence of client fingerprints within a visit points to the model
tier; and the agent's own answer to its user sometimes repeats the visited URL (in 5 of our 27 Claude
Code runs, all with WebFetch), so a per-visit identifier in that URL can reach the user's browser if
the user follows the link. This is UID smuggling~\cite{randall2022uid} with the agent as the carrier,
and it needs no cookie and no shared IP address. The per-visit sentence on our VM pages follows the
same logic: an ordinary sentence of the page, different on each visit, which ties a submitted summary
to the one visit an LLM read, even when the request is indistinguishable from a script that pauses.

\paragraph{Signals that did not work.} Several classic linking signals failed for coding agents. We
did not find TCP timestamps usable: current Linux randomises their offset per connection and Windows
does not send them by default. No agent connection resumed a TLS session, so session tickets link
nothing~\cite{sy2018tls}. In our local runs no agent requested \texttt{robots.txt}, \texttt{llms.txt} or a favicon
(the VM pages also link an \texttt{llms.txt} and Markdown copies; the campaign log tells whether any
agent follows them), and the cookie our server sets came back in none of 56 Claude Code requests, because each fetch or curl call
starts with an empty cookie jar. Cookies therefore link a person's browser visits but never an agent to
that browser, which is why linking has to rest on the network and the environment.

\section{Limitations}
\label{sec:limitations}

\paragraph{Scale.} Our study covers two agents, three members' machines (one Linux, two Windows) and two
sites, with tens of agent runs per condition. Several results rest on few runs: the model-tier pattern
on 12 to 19 runs per model, and the GPT side of the model classifier on 7 texts from one model. The
users are the three authors under fictitious personas, so linking is measured among three people; with
two Linux users or two identical Windows installs the environment signature would collide, and our
anonymity sets are an upper bound on how distinct users are.

\paragraph{Construction.} Our baseline scripts make the same requests as the task with no delay, so the
timing result separates an LLM in the loop from an automated script, not an agent from a script built to
pause; 48 of 48 deliberately slowed script runs were flagged as agents. The pause also measures
provider latency, which drifted by more than a second between two days. Our model classifier was trained
on output from the same CLIs the agents run, because we had no direct API access, so we cannot separate
the effect of the harness's system prompt from that of the agent session. Time between sessions is a
property of our campaign schedule as much as of the user, and we report linking without it.

\paragraph{Versions.} Agents update themselves often; we ran four Claude Code versions (2.1.319 to 2.1.332), and
one member's Codex CLI updated itself on a re-login during the study. Fingerprints of built-in tools, and which tool
a model chooses, can change with any release, so our numbers describe the versions we measured.

\section{Defenses}
\label{sec:defenses}

\paragraph{For agent developers.} The cheapest defense removes the per-visit identifier from what the
agent shows its user, or at least strips query strings and path tokens from cited URLs; that closes the
relay channel we observed. A built-in fetch tool with a common fingerprint (a widespread browser's TLS
and header profile) would remove the agent-specific signature, but routing all fetches through the
provider's servers, as chat assistants do, is the only change that removes the user's IP address and
environment as well, at the cost of local network access. Shelling out to the system curl hides the
agent among scripts but exposes the user's environment; agents that do so could use a bundled client
with a fixed, common fingerprint instead.

\paragraph{For users.} Running agents behind a VPN or proxy shared by many users, and in a container
with a common runtime, removes the IP and most of the environment signature; it does not remove the
timing pattern of an LLM in the loop or the agent's relay of URLs.

\paragraph{For sites that want to detect agents.} The same signals that track users can be used to
label agent traffic honestly. Here the defense and the attack coincide: a site that wants to treat
agents differently can rely on the timing pattern and on fetch-tool fingerprints, and agent developers
who want agents to be recognised can adopt explicit signals such as signed agent requests rather than
leaving sites to infer them.
